\documentclass[12pt]{article}
\usepackage[T1]{fontenc}
\usepackage[utf8]{inputenc}
\usepackage{lmodern}
\usepackage{textcomp}
\usepackage{enumitem}
\usepackage{geometry}
\usepackage{graphicx}
\usepackage{amsmath, amssymb}
\geometry{margin=1in}
\begin{document}
\begin{center}
History - Graduate Level Assessment 25 \
Total Marks: 40 \
Answer all questions.
\end{center}

\section*{Multiple Choice (10 marks)}
\begin{enumerate}[label=\arabic*.]
\item Which area is home to the tradition of making stone tools from chipped pebbles, called Hoabinhian?
\begin{enumerate}[label=(\alph*)]
    \item North America
    \item Australia
    \item Central America
    \item Arctic Region
    \item Middle East
\end{enumerate}
\item Which culture, previously-known as the temple-mound builders, based their subsistence primarily on maize and squash agriculture?
\begin{enumerate}[label=(\alph*)]
    \item Hohokam
    \item Folsom
    \item Adena
    \item Ancestral Puebloan
    \item Mississippian
\end{enumerate}
\item A(n) \_\_\_\_\_\_\_\_\_ contains primary refuse; a \_\_\_\_\_\_\_\_\_ contains secondary refuse.
\begin{enumerate}[label=(\alph*)]
    \item office; landfill
    \item dining area; garbage disposal
    \item bathroom; incinerator
    \item kitchen; garden
    \item playground; waste treatment plant
\end{enumerate}
\item Taphonomy is:
\begin{enumerate}[label=(\alph*)]
    \item the study of plant and animal life from a particular period in history.
    \item a study of the spatial relationships among artifacts, ecofacts, and features.
    \item the study of human behavior and culture.
    \item a technique to determine the age of artifacts based on their position in the soil.
    \item a dating technique based on the decay of a radioactive isotope of bone.
\end{enumerate}
\item Which culture lived a primarily sedentary lifestyle in the Southwest?
\begin{enumerate}[label=(\alph*)]
    \item Mogollon
    \item Navajo
    \item Clovis
    \item Paleo-Indians
    \item Hopewell
\end{enumerate}
\end{enumerate}

\section*{True / False (10 marks)}
\textit{Answer True or False}
\begin{enumerate}[label=\arabic*.]
\setcounter{enumi}{5}
\item True or False: This difference was particularly apparent in textual and visual works of early European studies of the Orient that positioned the East as irrational and backward in opposition to the rational and progressive West.
\item True or False: The Court of Justice held that Mr Steymann was entitled a different answer.
\item True or False: Although anthropologists worldwide refer to Tylor's definition of culture, in the a different answer "culture" emerged as the central and unifying concept of American anthropology, where it most commonly refers to the universal human capacity to classify and encode human experiences symbolically, and to communicate symbolically encoded experiences socially.
\item True or False: The military system became an intricate organization with the advance of the Empire. The Ottoman military was a complex system of recruiting and fief-holding. The main corps of the Ottoman Army included Janissary, Sipahi, Akıncı and Mehterân.
\item True or False: At Maecenas' insistence (according to the tradition) Virgil spent the ensuing years (perhaps 37-29 BC) on the long didactic hexameter poem called the Georgics (from Greek, "On Working the Earth") which he dedicated to Maecenas.
\end{enumerate}

\section*{Long Form Response (20 marks)}
\textit{Answer each question with a clear and complete explanation.}
\begin{enumerate}[label=\arabic*.]
\setcounter{enumi}{10}
\item Read the following passage and answer the question: One of the first recorded instances of translation in the West was the rendering of the Old Testament into Greek in the 3 rd century BCE. The translation is known as the "Septuagint", a name that refers to the seventy translators (seventy-two, in some versions) who were commissioned to translate the Bible at Alexandria, Egypt. Each translator worked in solitary confinement in his own cell, and according to legend all seventy versions proved identical. The Septuagint became the source text for later translations into many languages, including Latin, Coptic, Armenian and Georgian. Question: Why is the translation of the Old Testament into Greek known as the Septuagint?
\item Read the following passage and answer the question: Iranian art encompasses many disciplines, including architecture, painting, weaving, pottery, calligraphy, metalworking, and stonemasonry. The Median and Achaemenid empires left a significant classical art scene which remained as basic influences for the art of the later eras. Art of the Parthians was a mixture of Iranian and Hellenistic artworks, with their main motifs being scenes of royal hunting expeditions and investitures. The Sassanid art played a prominent role in the formation of both European and Asian medieval art, which carried forward to the Islamic world, and much of what later became known as Islamic learning, such as philology, literature, jurisprudence, philosophy, medicine, architecture, and science, were of Sassanid basis. Question: A blend of Iranian and what other type of artwork were the Parthians' art comprised of?
\end{enumerate}

\vfill
\noindent\textit{End of Paper}
\end{document}